\section{Supplement to Section 5: certificates}\label{supp:cert}

This section gives the details of the certificates of Section~\ref{M-sec:certificates} that are not needed to follow the argument: the vertex conditions and the enclosure of the periodic certificate C3, the validation protocols, the effect of the validity bound, what C3 gains over the averaged certificate C2 and how it converges in the number of segments, the results at low radial immersion and the straight-tooth counterexample to averaging.

\subsection{Vertex conditions and enclosure of the periodic certificate}\label{supp:cert-vertex}

On segment $k$, with $W(s)$ and $Y(s)$ from Eq.~\eqref{M-eq:Wt}, the decrease condition~\eqref{M-eq:lmi-per} is imposed at the six vertices of Proposition~\ref{M-prop:c3}, $(s,\sigma)=(0,0)$, $(\tfrac12,0)$ or $(1,1)$ and $r=\pm\varepsilon_k$, as
\begin{equation}\label{eq:app-per}
\begin{bmatrix}\Delta_kW-\Lambda-\Lambda^\top-2\lambda W(s) & BS_k-Y(s)^\top\\ S_kB^\top-Y(s) & 2S_k\end{bmatrix}\succ0,
\end{equation}
where $\Lambda=(A_0+Bk^\top)\,W(s)+E_r\,\big[\nu_0W_k+\nu_1W_{k+1}\big]$, $\nu_0=(c_{0,k}+r)(1-s)+c_{1,k}(s-\sigma)$ and $\nu_1=(c_{0,k}+r)s+c_{1,k}\sigma$; these reproduce $aH(t)W(t)$ exactly when $\sigma=s^2$. Since $E_r=p_rq_r^\top$, with $p_r=e_2/m$ and $q_r=e_1-e_{4+N}$, has rank one, the part of the matrix of Eq.~\eqref{eq:app-per} that depends on $r$ is $-r(\bar p\bar v^\top+\bar v\bar p^\top)$, with $\bar p=[p_r^\top\ 0]^\top$ and $\bar v=[(W(s)q_r)^\top\ 0]^\top$. By Petersen's lemma, exact for a scalar uncertainty \citep{petersen1987stabilization}, the matrix is positive definite for all $\lvert r\rvert\le\varepsilon_k$ if and only if, for some $\mu>0$,
\begin{equation}\label{eq:app-merged}
\begin{bmatrix}M_0-\mu\bar p\bar p^\top & \varepsilon_k\bar v\\ \varepsilon_k\bar v^\top & \mu\end{bmatrix}\succ0,
\end{equation}
with $M_0$ the matrix at $r=0$. Condition~\eqref{eq:app-merged} merges the vertices $r=\pm\varepsilon_k$ and leaves three LMIs per segment; for balance, $p_r$ is divided and $q_r$ multiplied by $10\,\big(\lVert p_r\rVert/(\lVert q_r\rVert\varepsilon_k)\big)^{1/2}$. The independent check tests the six vertices separately, with a round-off scale that also includes $\big(\lVert A_0+Bk^\top\rVert+(\lvert\nu_0\rvert+\lvert\nu_1\rvert)\lVert E_r\rVert\big)\max(\lVert W_k\rVert,\lVert W_{k+1}\rVert)$ and $(\lVert W_k\rVert+\lVert W_{k+1}\rVert)/\tau_s$, to cover rounding in $\Delta_kW$ and in $\mathcal AW$.

The line $c_{0,k}+c_{1,k}s$ is fitted by least squares to $aH$ at the midpoints of 2048 equal cells of segment $k$, and $\varepsilon_k$ is the largest residual plus $(L+\lvert c_{1,k}\rvert/\tau_s)\,\tau_s/4096$, where $L\ge\lvert\mathrm d(aH)/\mathrm dt\rvert$ holds because each tooth contributes to $\mathrm d(aH)/\mathrm dt$ only through the two ends of its engaged interval in $z$.
% src: src/c3_periodic.py (enclosure: N_ENC 2048; aH_rate_bound; corners; lmi: merged block, P_SCALE 10; check); README §14.2, §14.9

\subsection{Validation protocols}\label{supp:cert-protocol}

The off-axis test used boundary points of the certified ellipsoid (for C3, of the ellipsoid at the phase of application) in 18 directions: both signs of the points that maximise $\tilde x$, $\dot{\tilde x}$, $\tilde q$, $\dot{\tilde q}$, $\tilde x(t-T)$ and $\tilde x-\tilde x(t-T)$ on it, and six random points. They were applied in the simulator as full-state deviations at four phases a quarter of a tooth period apart, with the protocol of \ref{M-app:lmi}. These directions were sampled, not searched for the worst case.

The sampled check followed 48 trajectories per point of the collocated model with the true dead zone over 40 tooth periods, started at 0.999 of the boundary in random directions and phases. $\mathcal V$ never exceeded its initial value 0.998 for C2 and stayed below 0.997 after the first step for C3 (2256 starts). This is only a consistency check, as ellipsoids that are too large can pass it; the certificates rest on the independent check of Section~\ref{M-sec:cert-check}.

Floquet limits of the periodic collocated model used 200 Magnus steps per period and bisection to a relative $10^{-5}$. With $N+8$ nodes the spectral radius at the computed limit differs from 1 by at most $6\times10^{-6}$, and at nine speeds at half immersion the limit agrees with stored semi-discretisation limits \citep{insperger2002semi,insperger2004updated} within $1.3\times10^{-5}$; no such check was made at $a_e/D=0.1$.
% src: src/c2_certificate.py (boundary_directions, deviation_to_pert, td_offaxis, proof_check: nruns 48, nper 40); src/c3_periodic.py (td_offaxis, floquet_a_lim); results/c2_validate.npz (proof_vmax 0.998001, offaxis 0/3456); results/c3_summary.json (basin: proof_starts 2256, proof_vmax 0.9969, offaxis 0/3384); results/c3_floquet.npz (rho_n8: 4.9e-6, 5.5e-6); README §14.4 (SDM agreement 1.3e-5 at nine speeds)

\subsection{Validity bound}\label{supp:cert-delta}

Without the validity bound~\eqref{M-eq:lmi-valid}, the certified ellipsoid at \SI{2500}{\rpm} and $0.7\alim$ admitted histories displaced by 7--12 feeds, with the tool entirely out of the cut. These states return to the orbit in the linear model, but 32 of 72 off-axis runs from its boundary ended in chatter in the simulator, against none with the bound. Table~\ref{tab:delta} gives the certified kick for other values of $\delta$ and without the bound, with the command bound of the original C2 solves. Near the limit the sector condition sets the certificate and $\delta$ hardly matters; at shallow depths the bound sets it, and without the bound $\beta$ doubles at \SI{1090}{\rpm} and $0.15\alim$, but the set may then contain states with the tool out of the cut, as above.
% src: README §12.4 (7-12 f_z; 32 of 72 runs, 0 of 72 with the bound); verification/paper/revision_runs.json (beta_vs_delta; 23.367/11.689 = 2.0 at 1090 rpm, 0.15 a_lim; delta = 0.5 values equal results/c2_validate.npz cert)

\begin{table}[!tbp]
\centering
\caption{Certified kick $\beta$ (\si{\milli\metre\per\second}) of the averaged model (C2) against the validity bound $\delta$ of Eq.~\eqref{M-eq:lmi-valid}; $\delta=0.5$ is used in the article.}
\label{tab:delta}
\small
\begin{tabular}{@{}lcccc@{}}
\toprule
Point & $\delta=0.25$ & $\delta=0.5$ & $\delta=1$ & no bound\\
\midrule
\SI{1090}{\rpm}, $0.15\alim$ & 6.57 & 11.7 & 18.5 & 23.4\\
\SI{2870}{\rpm}, $0.3\alim$ & 5.69 & 8.74 & 10.0 & 10.4\\
\SI{2500}{\rpm}, $0.9\alim$ & 2.90 & 3.06 & 3.07 & 3.07\\
\bottomrule
\end{tabular}
\end{table}
% src: verification/paper/revision_runs.json (beta_vs_delta: 6.5748/11.6891/18.4958/23.3672; 5.6854/8.7429/10.0125/10.3694; 2.8990/3.0579/3.0712/3.0731 mm/s)

\subsection{What the periodic certificate gains, and its convergence}\label{supp:cert-decomp}

At $0.9\alim$ and the four speeds, relative to C2, the periodic coefficient with the single level of C2 gives 0.94--0.98, the averaged model with the segment levels 1.06--1.15, and both together (C3) 1.02--1.12: the periodic coefficient costs 2--6\,\%, and the segment levels gain 6--15\,\%. For the averaged, time-invariant model with a single level, a periodic $W$ cannot gain: the LMIs are then linear in $(W,Y,S)$ with constant coefficients and $\dot W$ integrates to zero over the period, so the period average of a feasible periodic certificate is a feasible constant one of the same size; on the periodic model the argument fails, since the period average of $aH(t)W(t)$ is not the product of the averages. C3 is not converged in the number of segments $M$: relative to C2, $M=8$ gives 0.93--1.03 and $M=32$ gives 1.08--1.22 at $0.9\alim$ (3--12\,\% above $M=16$), at about twice the memory. At $0.98\alim$, $M=32$ gives 5.15, 2.31 and \SI{5.01}{\milli\metre\per\second} at 1090, 1420 and \SI{2870}{\rpm}, against 4.84, 2.08 and \SI{4.83}{\milli\metre\per\second} for C2 and 4.60, 2.05 and \SI{1.70}{\milli\metre\per\second} for $M=16$. At \SI{2500}{\rpm} and $0.98\alim$ no candidate passed with either $M$: the best ones had a positive smallest eigenvalue below the round-off bound.
% src: README §14.5, §14.6; results/c3_decomp.npz, c3_edge.npz, c3_summary.json (decomp: averaged resolved 1.056-1.153, periodic single level 0.939-0.983, m_c 8 0.927-1.034, m_c 16 1.017-1.121, m_c 32 1.077-1.215; edge: c2 4.843/2.081/1.139/4.826, m_c_16 4.601/2.046/0/1.700, m_c_32 5.152/2.313/0/5.013 mm/s)

\subsection{Low radial immersion}\label{supp:cert-lowimm}

At $a_e/D=0.1$, $aH(t)$ varies more: over the four speeds it spans 0.66--1.40 times its mean at $0.5\alim$ and 0.79--1.24 times at $0.9\alim$, against 0.77--1.27 and 0.88--1.11 at half immersion, and the limits are 30--\SI{65}{\milli\metre} (14--\SI{30}{\milli\metre} at half immersion). All eight points (0.5 and $0.9\alim$) were certified, with C3/C2 = 0.94--1.04; the periodic and averaged critical kicks agree within 0.974--1.052, the certificate lies below $\Delta v^*$ in all 16 cases tested (four phases at the four points with finite $\Delta v^*$), and none of the 576 off-axis runs ended in chatter.
% src: README §14.4, §14.7; results/c3_lowimm.npz, c3_summary.json (lowimm: c3_over_c2 0.944-1.039, periodic_over_averaged_model 0.9739-1.0516, phase_cert_below_td 16/16, offaxis 0/576); aH ranges re-evaluated in closed form (0.660-1.403, 0.787-1.241; 0.766-1.272, 0.881-1.109)

\subsection{Straight teeth: a counterexample to averaging}\label{supp:cert-straight}

Averaging is known to fail in strongly interrupted cutting \citep{davies2002stability,bayly2003stability,merdol2004multi}. With straight teeth $aH(t)$ jumps as teeth enter and leave the cut, and at half immersion between 1600 and \SI{2100}{\rpm} (steps of \SI{25}{\rpm}) the uncontrolled cut has flip lobes below the averaged limit at 12 of 21 speeds (1675--\SI{1950}{\rpm}, from 0.947 of that limit). Where they persist with DVF ($g=\SI{100}{\volt\second\per\metre}$ at four speeds; the gain of the map at \SI{1950}{\rpm}), the forced orbit alone needs 27--\SI{66}{\volt} at 0.98 of the averaged limit, so $u_0<0$ and C2 claims nothing; gains of 250 and \SI{500}{\volt\second\per\metre} removed them, while gains of 141--\SI{250}{\volt\second\per\metre} and straight teeth at low immersion were not examined. Without control there is no forced command and the method of C2 does certify, which gives the counterexample of Table~\ref{tab:counter}: the averaged model, and with it the C2 certificate, is the same for both cutters, whereas the periodic model separates them. With straight teeth the cut is linearly unstable, though nearly benign in practice, its period-two oscillation staying at about \SI{1}{\micro\metre}.
% src: README §14.8; results/c3_straight.npz, c3_counter.npz (x_tooth alternation about 0.53 um, i.e. about 1 um peak to peak), c3_summary.json (straight, counter)

\begin{table}[!tbp]
\centering
\caption{Counterexample to averaging: \SI{1850}{\rpm}, no control ($g=0$), depth 0.98 times the averaged limit (\SI{23.6}{\milli\metre}), half immersion.}
\label{tab:counter}
\small
\begin{tabular}{@{}>{\raggedright\arraybackslash}p{0.34\textwidth}>{\raggedright\arraybackslash}p{0.29\textwidth}>{\raggedright\arraybackslash}p{0.27\textwidth}@{}}
\toprule
 & Straight teeth & Helical teeth (\ang{45})\\
\midrule
Critical Floquet multiplier & modulus 1.034 at \ang{180} & modulus 0.977 at \ang{177}\\
C2 certificate (averaged model) & \SI{13.2}{\milli\metre\per\second} & \SI{13.2}{\milli\metre\per\second}\\
C3 certificate (periodic model) & none (linearly unstable) & \SI{13.1}{\milli\metre\per\second}\\
Simulator after a \SI{0.1}{\milli\metre\per\second} kick & never settles; period-two oscillation of about \SI{1}{\micro\metre} & settles\\
\bottomrule
\end{tabular}
\end{table}
% src: README §14.8; results/c3_counter.npz (a 23.63 mm, mu_abs 1.0342/0.9767, mu_arg 180/177 deg, c2_size 13.201, c3_size 0/13.090 mm/s, sim_stable False/True), log_c3_counter.txt
